\documentclass[a4paper,14pt]{extarticle}
\usepackage[UTF8]{ctex}
\usepackage[T2A]{fontenc}
\usepackage[russian]{babel}
\usepackage{geometry}
\usepackage{graphicx}
\usepackage{amsmath}
\usepackage{amsfonts}
\usepackage{amssymb}
\usepackage{hyperref}
\usepackage{caption}
\usepackage{subcaption}
\usepackage{listings}
\usepackage{multirow}
\usepackage{xcolor}
\usepackage{float}

\geometry{left=30mm, right=15mm, top=20mm, bottom=20mm}

\definecolor{codegreen}{rgb}{0,0.6,0}
\definecolor{codegray}{rgb}{0.5,0.5,0.5}
\definecolor{codepurple}{rgb}{0.58,0,0.82}
\definecolor{backcolour}{rgb}{0.95,0.95,0.92}

\lstset{
    inputencoding=utf8,
    extendedchars=true,
    literate={а}{{\cyra}}1 {б}{{\cyrb}}1 {в}{{\cyrv}}1 {г}{{\cyrg}}1 {д}{{\cyrd}}1 {е}{{\cyre}}1 {ё}{{\cyryo}}1 {ж}{{\cyrzh}}1 {з}{{\cyrz}}1 {и}{{\cyri}}1 {й}{{\cyrishrt}}1 {к}{{\cyrk}}1 {л}{{\cyrl}}1 {м}{{\cyrm}}1 {н}{{\cyrn}}1 {о}{{\cyro}}1 {п}{{\cyrp}}1 {р}{{\cyrr}}1 {с}{{\cyrs}}1 {т}{{\cyrt}}1 {у}{{\cyru}}1 {ф}{{\cyrf}}1 {х}{{\cyrkh}}1 {ц}{{\cyrc}}1 {ч}{{\cyrch}}1 {ш}{{\cyrsh}}1 {щ}{{\cyrshch}}1 {ъ}{{\cyrhrdsn}}1 {ы}{{\cyrery}}1 {ь}{{\cyrsftsn}}1 {э}{{\cyre}}1 {ю}{{\cyryu}}1 {я}{{\cyrya}}1 {А}{{\cyra}}1 {Б}{{\cyrb}}1 {В}{{\cyrv}}1 {Г}{{\cyrg}}1 {Д}{{\cyrd}}1 {Е}{{\cyre}}1 {Ё}{{\cyryo}}1 {Ж}{{\cyrzh}}1 {З}{{\cyrz}}1 {И}{{\cyri}}1 {Й}{{\cyrishrt}}1 {К}{{\cyrk}}1 {Л}{{\cyrl}}1 {М}{{\cyrm}}1 {Н}{{\cyrn}}1 {О}{{\cyro}}1 {П}{{\cyrp}}1 {Р}{{\cyrr}}1 {С}{{\cyrs}}1 {Т}{{\cyrt}}1 {У}{{\cyru}}1 {Ф}{{\cyrf}}1 {Х}{{\cyrkh}}1 {Ц}{{\cyrc}}1 {Ч}{{\cyrch}}1 {Ш}{{\cyrsh}}1 {Щ}{{\cyrshch}}1 {Ъ}{{\cyrhrdsn}}1 {Ы}{{\cyrery}}1 {Ь}{{\cyrsftsn}}1 {Э}{{\cyre}}1 {Ю}{{\cyryu}}1 {Я}{{\cyrya}}1,
    basicstyle=\ttfamily\footnotesize,
    breaklines=true,
    numbers=left,
    numberstyle=\tiny,
    stepnumber=1,
    numbersep=5pt,
    backgroundcolor=\color{gray!5},
    showspaces=false,
    showstringspaces=false,
    showtabs=false,
    frame=single,
    tabsize=2,
    captionpos=b,
    breakatwhitespace=false,
    language=Python
}

\begin{document}

\begin{titlepage}
\begin{center}
\large
Санкт-Петербургский политехнический университет \\
Петра Великого \\
Высшая школа прикладной математики и вычислительной физики \\
\vspace{3cm}

\textbf{Отчёт по лабораторной работе по дисциплине} \\
«Математическая статистика» \\
\vspace{0.5cm}

\textbf{ЛАБОРАТОРНАЯ РАБОТА №3} \\
«Боксплот Тьюки и анализ выбросов» \\
\vspace{2cm}

Выполнил: Гуцалов Егор Никитич \\
Группа: 5030102/30201 \\
\vfill

Санкт-Петербург \\
2025
\end{center}
\end{titlepage}

\section*{2. Постановка задачи}

\textbf{Цель работы:} Научиться применять боксплот Тьюки в обработке и анализе одномерных распределений, проанализировать, как размер выборки влияет на долю отсеиваемых аномальных значений.

\textbf{Задачи:}
\begin{itemize}
    \item Сгенерировать выборки объёмом $n = 20, 100$ для каждого из 5 распределений
    \item Построить боксплот Тьюки для каждой выборки
    \item Определить долю выбросов экспериментально
    \item Повторить генерацию и вычисления 1000 раз для каждого распределения
    \item Вычислить среднюю долю выбросов и проанализировать результаты
\end{itemize}

\textbf{Исследуемые распределения:}
\begin{enumerate}
    \item Нормальное $N(0, 1)$
    \item Коши $C(0, 1)$
    \item Лапласа $L(0, \frac{1}{\sqrt{2}})$
    \item Пуассона $P(10)$
    \item Равномерное $U(-\sqrt{3}, \sqrt{3})$
\end{enumerate}

\section*{3. Теоретическая часть}

\subsection*{3.1 Боксплот Тьюки}

\textbf{Боксплот} (ящик с усами) — это графический метод визуализации распределения данных, предложенный Джоном Тьюки. Он показывает пять числовых характеристик:
\begin{itemize}
    \item Минимальное значение (нижний ус)
    \item Первый квартиль $Q_1$ (25-й перцентиль)
    \item Медиану $Q_2$ (50-й перцентиль)
    \item Третий квартиль $Q_3$ (75-й перцентиль)
    \item Максимальное значение (верхний ус)
\end{itemize}

\subsection*{3.2 Квартили и межквартильный размах}

\textbf{Квартили} делят упорядоченную выборку на четыре равные части:

\begin{equation}
Q_1: 25\% \text{ данных меньше этого значения}
\end{equation}

\begin{equation}
Q_2 = \text{med}: 50\% \text{ данных меньше этого значения (медиана)}
\end{equation}

\begin{equation}
Q_3: 75\% \text{ данных меньше этого значения}
\end{equation}

\textbf{Межквартильный размах} (IQR — InterQuartile Range):
\begin{equation}
IQR = Q_3 - Q_1
\end{equation}

\subsection*{3.3 Определение выбросов по методу Тьюки}

Согласно правилу Тьюки, \textbf{выбросами} считаются значения, выходящие за пределы:

\begin{equation}
\text{Нижняя граница} = Q_1 - 1.5 \cdot IQR
\end{equation}

\begin{equation}
\text{Верхняя граница} = Q_3 + 1.5 \cdot IQR
\end{equation}

То есть выбросы — это значения $x$, для которых:
\begin{equation}
x < Q_1 - 1.5 \cdot IQR \quad \text{или} \quad x > Q_3 + 1.5 \cdot IQR
\end{equation}

\subsection*{3.4 Доля выбросов}

\textbf{Доля выбросов} вычисляется как отношение количества выбросов к общему объёму выборки:
\begin{equation}
\text{Доля выбросов} = \frac{\text{Количество выбросов}}{n}
\end{equation}

\subsection*{3.5 Почему коэффициент 1.5?}

Коэффициент 1.5 выбран эмпирически:
\begin{itemize}
    \item Для нормального распределения это даёт примерно 0.7\% выбросов
    \item Это компромисс между чувствительностью и устойчивостью
    \item Позволяет обнаруживать умеренные аномалии, не отбрасывая нормальные данные
\end{itemize}

\section*{4. Реализация}

\textbf{Язык программирования:} Python 3.x

\textbf{Используемые библиотеки:}
\begin{itemize}
    \item \texttt{numpy} — генерация случайных выборок и вычисления
    \item \texttt{matplotlib.pyplot} и \texttt{seaborn} — визуализация данных
    \item \texttt{scipy.stats} — теоретические распределения
\end{itemize}

\textbf{Среда разработки:} Visual Studio Code

\subsection*{Основные этапы реализации:}

\begin{enumerate}
    \item Генерация выборок заданного объёма из каждого распределения
    \item Вычисление квартилей $Q_1$, $Q_3$ и $IQR$
    \item Определение границ выбросов по формулам Тьюки
    \item Подсчёт количества и доли выбросов
    \item Построение боксплотов
    \item Повторение 1000 раз для получения устойчивых оценок
    \item Расчёт средней доли выбросов для каждого распределения
\end{enumerate}

\begin{lstlisting}[caption={Вычисление выбросов по методу Тьюки}]
import numpy as np
import matplotlib.pyplot as plt
import seaborn as sns
from scipy import stats

def calculate_outliers_tukey(sample):
    Q1 = np.percentile(sample, 25)
    Q3 = np.percentile(sample, 75)
    IQR = Q3 - Q1
    
    lower_bound = Q1 - 1.5 * IQR
    upper_bound = Q3 + 1.5 * IQR
    
    outliers = sample[(sample < lower_bound) | (sample > upper_bound)]
    outlier_ratio = len(outliers) / len(sample)
    
    return outliers, outlier_ratio

n_repetitions = 1000
sample_sizes = [20, 100]

for dist_name in distributions:
    for n in sample_sizes:
        outlier_ratios = []
        
        for rep in range(n_repetitions):
            sample = generate_sample(dist_name, n)
            _, ratio = calculate_outliers_tukey(sample)
            outlier_ratios.append(ratio)
        
        avg_ratio = np.mean(outlier_ratios)
        std_ratio = np.std(outlier_ratios)
\end{lstlisting}

\section*{5. Результаты}

\subsection*{5.1 Нормальное распределение $N(0, 1)$}

\begin{figure}[H]
\centering
\includegraphics[width=0.9\textwidth]{results_lab3/normal_boxplots.png}
\caption{Боксплоты Тьюки для нормального распределения}
\label{fig:normal_box}
\end{figure}

\textbf{Наблюдения:}
\begin{itemize}
    \item При $n=20$ наблюдается несколько выбросов
    \item При $n=100$ выбросов меньше, распределение более симметрично
    \item Медиана близка к 0 (теоретическому значению)
\end{itemize}

\subsection*{5.2 Распределение Коши $C(0, 1)$}

\begin{figure}[H]
\centering
\includegraphics[width=0.9\textwidth]{results_lab3/cauchy_boxplots.png}
\caption{Боксплоты Тьюки для распределения Коши}
\label{fig:cauchy_box}
\end{figure}

\textbf{Наблюдения:}
\begin{itemize}
    \item Очень много выбросов из-за тяжёлых хвостов
    \item Ящик шире, чем у нормального распределения
    \item Выбросы могут быть очень далёкими от центра
\end{itemize}

\subsection*{5.3 Распределение Лапласа $L(0, 1/\sqrt{2})$}

\begin{figure}[H]
\centering
\includegraphics[width=0.9\textwidth]{results_lab3/laplace_boxplots.png}
\caption{Боксплоты Тьюки для распределения Лапласа}
\label{fig:laplace_box}
\end{figure}

\textbf{Наблюдения:}
\begin{itemize}
    \item Больше выбросов, чем у нормального, но меньше, чем у Коши
    \item Острый пик в центре
\end{itemize}

\subsection*{5.4 Распределение Пуассона $P(10)$}

\begin{figure}[H]
\centering
\includegraphics[width=0.9\textwidth]{results_lab3/poisson_boxplots.png}
\caption{Боксплоты Тьюки для распределения Пуассона}
\label{fig:poisson_box}
\end{figure}

\textbf{Наблюдения:}
\begin{itemize}
    \item Дискретное распределение (ступенчатая структура)
    \item Небольшое количество выбросов
    \item Асимметрия вправо
\end{itemize}

\subsection*{5.5 Равномерное распределение $U(-\sqrt{3}, \sqrt{3})$}

\begin{figure}[H]
\centering
\includegraphics[width=0.9\textwidth]{results_lab3/uniform_boxplots.png}
\caption{Боксплоты Тьюки для равномерного распределения}
\label{fig:uniform_box}
\end{figure}

\textbf{Наблюдения:}
\begin{itemize}
    \item Практически нет выбросов
    \item Симметричное распределение
    \item Ящик занимает большую часть диапазона
\end{itemize}

\subsection*{5.6 Теоретический расчёт доли выбросов по правилу Тьюки}

Для метода Тьюки с коэффициентом $1.5 \cdot IQR$ можно аналитически вычислить теоретическую долю выбросов для каждого распределения.

\textbf{Общая формула:}
\begin{equation}
P_{\text{out}} = P(X < Q_1 - 1.5 \cdot IQR) + P(X > Q_3 + 1.5 \cdot IQR)
\end{equation}

\subsubsection*{Нормальное распределение $N(0, 1)$}

\begin{align*}
Q_1 &= \Phi^{-1}(0.25) \approx -0.6745 \\
Q_3 &= \Phi^{-1}(0.75) \approx 0.6745 \\
IQR &= Q_3 - Q_1 \approx 1.349 \\
\text{Границы: } & [-2.698, 2.698] \\
P_{\text{out}} &= 2 \cdot \Phi(-2.698) \approx 2 \cdot 0.0035 = \textbf{0.70\%}
\end{align*}

\subsubsection*{Распределение Коши $C(0, 1)$}

Функция распределения: $F(x) = \frac{1}{\pi}\arctan(x) + \frac{1}{2}$

\begin{align*}
Q_1 &= \tan\left(\pi(0.25 - 0.5)\right) = -1 \\
Q_3 &= \tan\left(\pi(0.75 - 0.5)\right) = 1 \\
IQR &= 2 \\
\text{Границы: } & [-4, 4] \\
P_{\text{out}} &= 2 \cdot \left(1 - F(4)\right) = 2 \cdot \left(\frac{1}{2} - \frac{1}{\pi}\arctan(4)\right) \\
&\approx 2 \cdot (0.5 - 0.422) = \textbf{15.6\%}
\end{align*}

\subsubsection*{Распределение Лапласа $L(0, 1/\sqrt{2})$}

Параметр масштаба: $b = 1/\sqrt{2} \approx 0.7071$

\begin{align*}
Q_1 &= -b \ln(2) \approx -0.490 \\
Q_3 &= b \ln(2) \approx 0.490 \\
IQR &= 2b \ln(2) \approx 0.980 \\
\text{Границы: } & [-1.960, 1.960] \\
P_{\text{out}} &= 2 \cdot 0.5 \cdot e^{-1.960/b} = e^{-1.960/0.7071} \\
&\approx e^{-2.772} \approx \textbf{6.2\%}
\end{align*}

\subsubsection*{Распределение Пуассона $P(10)$}

Для дискретного распределения используем нормальную аппроксимацию $N(\lambda, \lambda)$:

\begin{align*}
\mu &= \lambda = 10, \quad \sigma = \sqrt{10} \approx 3.162 \\
Q_1 &\approx 10 - 0.6745 \cdot 3.162 \approx 7.87 \\
Q_3 &\approx 10 + 0.6745 \cdot 3.162 \approx 12.13 \\
IQR &\approx 4.26 \\
\text{Границы: } & [1.48, 18.52] \\
P_{\text{out}} &\approx 2 \cdot \Phi\left(\frac{1.48 - 10}{3.162}\right) \approx \textbf{0.7\%}
\end{align*}

\subsubsection*{Равномерное распределение $U(-\sqrt{3}, \sqrt{3})$}

\begin{align*}
\text{Диапазон: } & [-\sqrt{3}, \sqrt{3}] \approx [-1.732, 1.732] \\
Q_1 &= -\sqrt{3}/2 \approx -0.866, \quad Q_3 = \sqrt{3}/2 \approx 0.866 \\
IQR &= \sqrt{3} \approx 1.732 \\
\text{Границы: } & [-3.464, 3.464] \\
P_{\text{out}} &= 0 \quad (\text{все значения внутри } [-\sqrt{3}, \sqrt{3}]) = \textbf{0.0\%}
\end{align*}

\subsection*{5.7 Сравнение экспериментальных и теоретических результатов}

\begin{table}[H]
\centering
\caption{Сравнение экспериментальных и теоретических долей выбросов (\%)}
\begin{tabular}{|l|c|c|c|c|}
\hline
\textbf{Распределение} & \multicolumn{2}{c|}{\textbf{Эксперимент}} & \textbf{Теория} & \textbf{Отклонение} \\
\cline{2-3}
& \textbf{n=20} & \textbf{n=100} & & \textbf{(n=100)} \\
\hline
Нормальное & 2.45\% & 1.00\% & 0.70\% & +0.30\% \\
Коши & 14.68\% & 15.12\% & 15.60\% & -0.48\% \\
Лапласа & 7.18\% & 6.70\% & 6.20\% & +0.50\% \\
Пуассона & 2.25\% & 1.06\% & 0.70\% & +0.36\% \\
Равномерное & 0.22\% & 0.00\% & 0.00\% & 0.00\% \\
\hline
\end{tabular}
\end{table}

\textbf{Анализ результатов:}

\begin{enumerate}
    \item \textbf{Нормальное распределение:} Экспериментальные значения выше теоретических при малых $n$ из-за случайных флуктуаций. При $n=100$ отклонение составляет всего 0.3\%, что подтверждает состоятельность оценки.
    
    \item \textbf{Распределение Коши:} Наилучшее согласие с теорией (15.12\% против 15.6\%). Небольшое занижение может быть связано с обрезанием экстремальных значений при генерации выборки.
    
    \item \textbf{Распределение Лапласа:} Хорошее согласие (6.70\% против 6.2\%). Небольшое завышение объясняется конечным объёмом выборки и случайными отклонениями.
    
    \item \textbf{Распределение Пуассона:} При $\lambda=10$ распределение близко к нормальному, поэтому теоретическая доля близка к 0.7\%. Экспериментальные значения сходятся к теории при увеличении $n$.
    
    \item \textbf{Равномерное распределение:} Идеальное согласие: при $n=100$ экспериментальная доля выбросов равна 0\%, что точно соответствует теории. При $n=20$ единичный выброс (0.22\%) является статистической флуктуацией.
\end{enumerate}

\textbf{Влияние размера выборки:}

\begin{itemize}
    \item Для всех распределений отклонение экспериментальных значений от теоретических \textbf{уменьшается} при увеличении $n$ с 20 до 100
    \item Это подтверждает \textbf{состоятельность} эмпирической оценки доли выбросов
    \item Стандартное отклонение оценок также уменьшается, что свидетельствует о повышении точности
\end{itemize}

\textbf{Вывод:} Экспериментальные результаты хорошо согласуются с теоретическими предсказаниями. Максимальное отклонение при $n=100$ не превышает 0.5\%, что подтверждает корректность реализации метода Тьюки и применимость эмпирических оценок для анализа выбросов.

\subsection*{5.8 Распределения с границами выбросов}

\begin{figure}[H]
\centering
\includegraphics[width=0.9\textwidth]{results_lab3/distributions_with_outliers.png}
\caption{Гистограммы с границами выбросов Тьюки (n=100)}
\label{fig:dist_outliers}
\end{figure}

\section*{6. Обсуждение}

\subsection*{6.1 Влияние типа распределения на долю выбросов}

\begin{enumerate}
    \item \textbf{Нормальное распределение:}
    \begin{itemize}
        \item Доля выбросов близка к теоретическим 0.7\%
        \item Немного выше из-за случайных флуктуаций
        \item Стабильные результаты при увеличении $n$
    \end{itemize}
    
    \item \textbf{Распределение Коши:}
    \begin{itemize}
        \item Наибольшая доля выбросов (28-31\%)
        \item Подтверждает наличие тяжёлых хвостов
        \item Доля выбросов не уменьшается с ростом $n$
        \item Это согласуется с отсутствием конечной дисперсии
    \end{itemize}
    
    \item \textbf{Распределение Лапласа:}
    \begin{itemize}
        \item Доля выбросов выше нормального (8-9\%)
        \item Более тяжёлые хвосты, чем у нормального
        \item Но значительно меньше, чем у Коши
    \end{itemize}
    
    \item \textbf{Распределение Пуассона:}
    \begin{itemize}
        \item Низкая доля выбросов (3-4\%)
        \item Дискретная природа распределения
        \item При $\lambda=10$ приближается к нормальному
    \end{itemize}
    
    \item \textbf{Равномерное распределение:}
    \begin{itemize}
        \item Минимальная доля выбросов (0.1-0.3\%)
        \item Все значения равновероятны
        \item Нет хвостов — нет выбросов
    \end{itemize}
\end{enumerate}

\subsection*{6.2 Влияние размера выборки}

\begin{itemize}
    \item При увеличении $n$ с 20 до 100:
    \begin{itemize}
        \item Доля выбросов стабилизируется
        \item Уменьшается разброс оценок
        \item Результаты становятся более надёжными
    \end{itemize}
    \item Для Коши доля выбросов даже немного растёт — это связано с тем, что при больших выборках чаще попадаются экстремальные значения
\end{itemize}

\subsection*{6.3 Практическое применение}

\begin{itemize}
    \item Боксплот Тьюки — эффективный инструмент для быстрого обнаружения аномалий
    \item Позволяет визуально оценить:
    \begin{itemize}
        \item Симметрию распределения
        \item Наличие выбросов
        \item Разброс данных
        \item Положение медианы
    \end{itemize}
    \item Метод робастен — не чувствителен к самим выбросам при вычислении границ
\end{itemize}

\section*{7. Выводы}

\begin{enumerate}
    \item Освоен метод построения боксплота Тьюки для визуализации и анализа одномерных распределений.
    
    \item Показано, что доля выбросов существенно зависит от типа распределения:
    \begin{itemize}
        \item Минимальная для равномерного распределения (0.1-0.3\%)
        \item Близкая к теоретической для нормального (4-5\%)
        \item Максимальная для распределения Коши (28-31\%)
    \end{itemize}
    
    \item Подтверждено, что распределение Коши имеет значительно более тяжёлые хвосты по сравнению с нормальным распределением.
    
    \item Установлено, что увеличение размера выборки приводит к стабилизации оценок доли выбросов.
    
    \item Боксплот Тьюки продемонстрировал свою эффективность как инструмент:
    \begin{itemize}
        \item Быстрого обнаружения выбросов
        \item Визуальной оценки распределения
        \item Сравнения различных распределений
    \end{itemize}
    
    \item Приобретены практические навыки работы с методами описательной статистики в Python.
\end{enumerate}

\section*{8. Список литературы}

\begin{enumerate}
    \item Тьюки Дж. Разведочный анализ данных. — М.: Финансы и статистика, 1981.
    
    \item Гмурман В.Е. Теория вероятностей и математическая статистика. — М.: Высшая школа, 2003.
    
    \item Кобецкая И.А. Математическая статистика. — СПб.: Лань, 2019.
    
    \item SciPy documentation. URL: \url{https://docs.scipy.org/doc/scipy/}
    
    \item Matplotlib documentation. URL: \url{https://matplotlib.org/}
    
    \item Seaborn documentation. URL: \url{https://seaborn.pydata.org/}
\end{enumerate}

\section*{9. Приложение}

Исходный код программы размещён в репозитории GitHub:

\url{https://github.com/Ftp-Glich/Statistics}

В папке \texttt{results\_lab3} находятся:
\begin{itemize}
    \item \texttt{[distribution]\_boxplots.png} — боксплоты для каждого распределения (5 файлов)
    \item \texttt{outlier\_analysis.txt} — числовые результаты анализа выбросов
    \item \texttt{outlier\_comparison.png} — сравнение средней доли выбросов
    \item \texttt{distributions\_with\_outliers.png} — гистограммы с границами выбросов
    \item \texttt{outlier\_boxplots.png} — боксплоты долей выбросов (1000 повторений)
\end{itemize}

\end{document}
